% Chapter 6 — Discussion

\chapter{Discussion}
\label{chap:discussion}

The experiments in Chapter~\ref{chap:results} show that the comparison between
the learners depends strongly on the regret notion and the experimental
setting. This chapter returns to the research questions from
Chapter~\ref{chap:introduction}: how the algorithms differ across the three
regret notions, how stable these differences are across experimental settings,
and how they appear in repeated games.

% ============================================================
\section{Regret Metrics and Theoretical Comparison}
\label{sec:discussion-theory-and-metrics}

The experiments measure expected action regret, while the guarantees in
Chapter~\ref{chap:algorithms} use several related regret notions. A few simple
relations are enough to interpret most of the results.

First, action regret \(A_T\) and distribution regret \(D_T\) differ through
the randomness of the sampled actions. In the bounded finite-action setting,
this sampling difference is of root-\(T\) order,
\[
    A_T-D_T=O_p(\sqrt{T}),
\]
up to logarithmic factors for the finite deviation class
\citep{greenwald_bounds_2006}. Thus a persistent super-root-\(T\) action-regret
rate cannot be explained by sampling alone and implies the same leading
polynomial order for distribution regret. Around the root-\(T\) scale, however,
the two cannot be separated reliably by the current evaluator.

Second, several bandit guarantees concern pseudo-regret
\[
    P_T=\max_F \mathbb{E}[Z_F],
\]
whereas the experiments estimate expected action regret
\[
    E_T=\mathbb{E}\!\left[\max_F Z_F\right].
\]
Since
\[
    P_T\le E_T,
\]
sublinear expected action regret also implies sublinear pseudo-regret, but large
expected action regret does not imply large pseudo-regret. This distinction is
made explicit by \citet{huang_near-optimal_2023}.

Some references give stronger probabilistic guarantees than the replicate
means measured here. Huang and Pan derive an expected swap-regret bound from
a high-probability instantaneous-regret bound
\citep{huang_near-optimal_2023}. Such high-probability statements are not
tested directly by the mean regret curves used in the experiments.

These relations give a more precise interpretation of the individual
algorithms.

\paragraph{Hedge.}
Its external action-regret exponents stay close to \(0.5\). Since action and
distribution regret can differ by root-\(T\) sampling fluctuations, this does
not determine the corresponding distribution-regret rate and therefore does
not directly test the standard Hedge guarantee \citep{freund_decision-theoretic_1997}.
More interestingly, in RPSLS self-play its internal and swap exponents rise
to about \(0.92\) and \(0.96\)--\(0.97\). If this trend persists, sampling
noise cannot explain the difference: the corresponding distribution regrets
must also grow faster than root-\(T\).

\paragraph{OptHedge.}
The realized internal- and swap-regret exponents are above the root scale in
several experiments. These quantities do not directly test the faster
standalone OptHedge result of \citet{chen_hedging_2020}, which concerns
external mixed regret in optimistic self-play. The theoretical sub-root rate
therefore remains unresolved by the current evaluator.

\paragraph{BM-Hedge.}
Its swap action-regret exponents are generally close to \(0.5\). Since this is
also the scale of the action--distribution sampling difference, the experiment
does not determine whether the corresponding distribution regret follows the
root-\(T\) Blum--Mansour bound or a smaller rate
\citep{blum_external_2007}.

\paragraph{BM-OptHedge.}
This method has some of the smallest internal- and swap-regret exponents in the
one-player experiments. The empirical result is favorable, but the
sub-root mixed-regret guarantee of \citet{chen_hedging_2020} cannot be
identified directly from sampled action regret.

\paragraph{Ito-Hedge.}
Ito's reduction gives a swap pseudo-regret guarantee
\citep{ito_tight_2020}. The observed expected swap regret is sublinear in the
tested settings, so \(P_T\le E_T\) also gives sublinear swap pseudo-regret if
the observed trend persists. The exact pseudo-regret rate cannot be recovered
from the current metric.

\paragraph{RM.}
RM shows sublinear internal and swap regret in the experiments, including its
RPSLS profiles. This is compatible with the convergence behavior proved for
the Hart--Mas-Colell repeated-game procedure
\citep{hart_simple_2000}. That theorem is a self-play result, however, so the
one-player and arbitrary cross-play experiments should mainly be viewed as
empirical comparisons.

\paragraph{SRM.}
SRM also shows sublinear expected internal and swap regret, often with exponents
near or below the root scale. This is consistent with its almost-sure
pairwise-regret result under the assumptions of
\citet{hart_simple_2000}, although expectation-level experiments cannot verify
almost-sure convergence.

\paragraph{EXP3.}
Its expected external action regret remains roughly at the root scale, while
internal and swap regret can grow substantially faster. Since the standard
guarantee concerns external pseudo-regret
\citep{auer_nonstochastic_2002}, the sublinear external results also support
sublinear pseudo-regret. The larger stronger-regret values do not contradict
the external-regret guarantee.

\paragraph{EXP3-IX.}
Neu proves a high-probability bound for realized external regret
\citep{neu_explore_2015}. The experiments instead report the mean realized
regret across replicates. Its fitted external-regret exponents remain close to
the root-$T$ scale, but this expectation-level result does not establish the
stronger high-probability guarantee.

\paragraph{BM-EXP3.}
Its expected swap regret is sublinear and usually close to the root scale.
Because the literature guarantee is on swap pseudo-regret
\citep{blum_external_2007}, this also supports sublinear pseudo-regret, although
the fitted action-regret exponent should not be interpreted as the
pseudo-regret exponent.

\paragraph{Ito-Tsallis.}
The same argument applies here. Expected swap regret remains sublinear in all
tested environments, so the weaker swap pseudo-regret is also sublinear if the
trend persists. The relatively larger exponent in HFA cannot be used to reject
a sharper pseudo-regret bound.

\paragraph{LCE-IX.}
This is the cleanest comparison because \citet{huang_near-optimal_2023}
directly bound expected swap action regret. The fitted swap-regret exponents
are close to \(0.5\) in both one-player environments, matching the root-\(T\)
dependence of the theoretical bound for fixed \(K\). The experiment still does
not test the stronger high-probability statement.

Overall, the empirical metric is informative in several directions. Large
super-root action regret can rule out sampling as the explanation for a
distributional effect, while small expected action regret can imply small
pseudo-regret. Other comparisons remain unresolved, especially sub-root
mixed-strategy rates and stronger probabilistic guarantees. These are
limitations of the current evaluator rather than contradictions with the
corresponding theory.

% ============================================================
\section{Robustness Across Experimental Conditions}
\label{sec:discussion-robustness}

The main empirical pattern is that the separation between algorithms is not
equally strong in every environment. HFA produces much clearer differences in
internal and swap regret than LRW, where many of the learners remain close
together.

A plausible reason is the feedback structure of the two reward processes. LRW
evolves independently of the learner, whereas HFA reacts to past action
frequencies. An update that changes the learner's behavior can therefore also
change its future rewards in HFA. This feedback loop gives small differences
between algorithms more opportunity to accumulate, which fits the stronger
separation observed in the experiments.

The action-space experiments show a related finite-horizon effect. Increasing
\(K\) generally raises swap regret and delays the main decline of the regret
curves. The horizon fits likewise vary across environments: the exponent
\(\alpha\) alone does not determine which algorithm looks better at a practical
horizon, since the coefficient \(c\) can shift the curves substantially.

The results therefore suggest a fairly simple picture. Algorithms designed for
stronger regret usually become more clearly separated from external-regret
baselines when internal or swap regret is measured, but the size of this
difference depends on the environment, action space, and horizon.

A broader difference is also visible between the two feedback settings.
Under full information, several swap-regret-oriented methods show clearer
separation from the external-regret baseline, and some achieve noticeably
smaller fitted internal- and swap-regret exponents. Under bandit feedback, the
corresponding differences are generally less pronounced, with most fitted
exponents remaining closer to the root-$T$ scale. This is consistent with the
additional uncertainty introduced by observing only the reward of the sampled
action. Since the two settings contain different algorithm families, however,
this should be read as a qualitative comparison rather than a controlled
feedback ablation.

% ============================================================
\section{Equilibrium and Opponent Effects in Repeated Games}
\label{sec:discussion-games}

Two patterns stand out in the repeated-game experiments.

First, the difference between regret notions is also visible in empirical
equilibrium behavior. In RPS, Hedge self-play gets much closer to the CCE set
than to the CE set, while several stronger-regret methods approach CE more
closely. This matches the standard link between external regret and CCE, and
between action-dependent regret and CE
\citep{hart_simple_2000,blum_external_2007}.

The joint-action distributions make this intuitive. Hedge can have nearly
uniform marginal action frequencies while still showing a structured joint
distribution, so balanced marginals alone do not remove profitable conditional
deviations.

The trajectories show a similar distinction. For Hedge, CE and CCE distance
first decrease together; later, CCE distance stays near zero while CE distance
rises again. A natural interpretation is that fixed deviations are already
unprofitable while some action-dependent deviations remain. A weaker version
of the same pattern appears for EXP3 and EXP3-IX.

Second, regret depends strongly on the opponent. Hedge has very large internal
and swap exponents in self-play, but both move much closer to the root scale
against BM-Hedge, Ito-Hedge, and SRM. EXP3 shows the same tendency more weakly
under bandit feedback.

The RM--SRM profiles give the clearest asymmetric example. SRM has internal
and swap exponents around \(0.50\) in SRM--SRM self-play, but they fall to
about \(0.34\) and \(0.36\) against RM, while RM does not improve to the same
extent. This suggests that the opponent can affect both regret levels and their
finite-horizon growth.

Taken together, the repeated-game results support the expected connection
between stronger regret and CE behavior, while showing that finite-horizon
dynamics depend strongly on the interaction between the two learners.
